% ========== Cover & Basic Information ==========
\title{ElegantBook Test Document}
\subtitle{Core Features of the Template}
\author{Test Engineer}
\institute{Template Testing Team}
\date{\today}
\version{v1.0}
\extrainfo{This document is only for template function testing and has no academic value}
\bioinfo{Test Batch}{TEST-FULL-20260506}

% ========== Table of Contents & Basic Settings ==========
\setcounter{tocdepth}{3} % Show table of contents up to level 3

% ========== Test References ==========
\addbibresource[location=local]{reference.bib}

\begin{document}

\maketitle
\frontmatter
\tableofcontents

% ========== Preface (with Section Subsections) ==========
\pagenumbering{Roman}
\setcounter{page}{0}
\chapter*{Test Preface}
\markboth{Test Preface}{}
\section*{Document Description}
\markright{Document Description}
This document is a full-function test case for ElegantBook, covering all basic and advanced features of the template. All content is fictional test data with no practical significance.
\newpage
\section*{Test Scope}
\markright{Test Scope}
This test includes: cover, table of contents, fonts, multi-level lists, complex formulas, all mathematical environments, cross-references, figures and tables, marginal notes, chapter summaries, exercises, references, and appendices.

\mainmatter

% ========== Chapter 1: Basic Typesetting & Multi-level Lists ==========
\chapter{Basic Typesetting and List Testing}
\begin{introduction}[Test Items of This Chapter]
\item Chinese Fonts and Text Formatting
\item 3-level Unordered List Nesting
\item 3-level Ordered List Nesting
\item Ordinary Paragraphs and Emphasized Text
\end{introduction}

\section{Font and Text Testing}
Normal body text test. \textbf{Bold}, \textit{Italic}, \underline{Underline}, \textcolor{red}{Red Text Test}.

\section{3-level Unordered List Test}
\begin{itemize}
\item Level 1 Test Item A
\item Level 1 Test Item B
    \begin{itemize}
    \item Level 2 Test Item B1
    \item Level 2 Test Item B2
        \begin{itemize}
        \item Level 3 Test Item B2-1
        \item Level 3 Test Item B2-2
            \begin{itemize}
            \item Level 4 Test Item B2-1a
            \item Level 4 Test Item B2-2b
            \end{itemize}
        \end{itemize}
    \end{itemize}
\end{itemize}

\section{3-level Ordered List Test}
\begin{enumerate}
\item Level 1 Ordered Test 1
\item Level 1 Ordered Test 2
    \begin{enumerate}
    \item Level 2 Ordered Test 2-1
    \item Level 2 Ordered Test 2-2
        \begin{enumerate}
        \item Level 3 Ordered Test 2-2-1
        \item Level 3 Ordered Test 2-2-2
            \begin{enumerate}
            \item Level 4 Ordered Test 2-2-2-1
            \item Level 4 Ordered Test 2-2-2-2
            \end{enumerate}
        \end{enumerate}
    \end{enumerate}
\end{enumerate}

% ========== Chapter 2: Complex Mathematical Formula Testing ==========
\chapter{Complex Mathematical Formula Testing}
\begin{introduction}
\item Single-line Complex Formula
\item Multi-line Formula (align)
\item Matrix/Determinant
\item Multiple Integrals/Summations/Limits/Fractions
\item Formula Labels and Cross-references
\end{introduction}

\section{Single-line Complex Formula}
Inline complex formulas: $\sum_{i=1}^n \frac{x_i^2}{\sqrt{y_i+1}} \geq \bar{x}^2$, $\lim_{n\to\infty} \left(1+\frac{1}{n}\right)^n = e$.

\section{Multi-line Formulas and Matrices}
\begin{align}
a^2 + b^2 &= c^2 \label{eq:gougu} \\
\int_{0}^{1}\int_{0}^{1} f(x,y) dxdy &= 1 \label{eq:double} \\
\begin{vmatrix}
a & b \\
c & d
\end{vmatrix} &= ad-bc \label{eq:det}
\end{align}

\section{Large Operator Formulas}
\begin{equation}\label{eq:series}
\sum_{k=0}^{\infty} \frac{x^k}{k!} = e^x,\quad \oint_{C} Pdx+Qdy = \iint_{D}\left(\frac{\partial Q}{\partial x}-\frac{\partial P}{\partial y}\right)dxdy
\end{equation}

Formula cross-reference: Pythagorean theorem \eqref{eq:gougu}, double integral \eqref{eq:double}, determinant \eqref{eq:det}, series \eqref{eq:series}.

% ========== Chapter 3: Full Mathematical Environment Testing ==========
\chapter{Full Mathematical Environment Testing}
\begin{introduction}
\item Definition/Theorem/Lemma/Corollary/Axiom/Postulate
\item Proposition/Example/Problem/Exercise
\item Note/Remark/Conclusion/Assumption/Property
\item Solution/Proof Environment
\end{introduction}

\section{Environment Testing}

% 1. Definition
\ifdefined\simpleTest
\begin{definition}[Test Definition]\label{def:test}
\else
\begin{definition}[Test Definition]{test}
\fi
This is the test content of the definition environment, with no practical mathematical meaning, only for style verification.
\end{definition}

% 2. Theorem
\ifdefined\simpleTest
\begin{theorem}[Test Theorem]\label{thm:test}
\else
\begin{theorem}[Test Theorem]{test}
\fi
If $x>0$, then $x+1>1$, corresponding to Formula \eqref{eq:gougu}.
\end{theorem}

% 3. Lemma
\ifdefined\simpleTest
\begin{lemma}[Test Lemma]\label{lem:test}
\else
\begin{lemma}[Test Lemma]{test}
\fi
For any real number $x$, $x^2\geq0$ holds.
\end{lemma}

% 4. Corollary
\ifdefined\simpleTest
\begin{corollary}[Test Corollary]\label{cor:test}
\else
\begin{corollary}[Test Corollary]{test}
\fi
From Lemma \ref{lem:test}, we can get: $(x-1)^2\geq0$.
\end{corollary}

% 5. Axiom
\ifdefined\simpleTest
\begin{axiom}[Test Axiom]\label{axi:test}
\else
\begin{axiom}[Test Axiom]{test}
\fi
Test axiom, no practical meaning.
\end{axiom}

% 6. Postulate
\ifdefined\simpleTest
\begin{postulate}[Test Postulate]\label{pos:test}
\else
\begin{postulate}[Test Postulate]{test}
\fi
Test postulate, only for environment verification.
\end{postulate}

% 7. Proposition
\ifdefined\simpleTest
\begin{proposition}[Test Proposition]\label{pro:test}
\else
\begin{proposition}[Test Proposition]{test}
\fi
Test proposition content, associated with Definition \ref{def:test}.
\end{proposition}

% 8. Example
\begin{example}\label{exa:test}
Example environment test, automatic numbering.
\end{example}

% 9. Problem
\begin{problem}\label{probl:test}
Problem environment test, used to propose unsolved problems.
\end{problem}

% 10. Exercise
\begin{exercise}\label{exer:test}
Calculate: $1+2+3+\dots+10$.
\end{exercise}

% 11. Note
\begin{note}
Note environment: no numbering, for key reminders.
\end{note}

% 12. Remark
\begin{remark}
Remark environment: for supplementary explanations.
\end{remark}

% 13. Conclusion
\begin{conclusion}
Conclusion environment: test summary content.
\end{conclusion}

% 14. Assumption
\begin{assumption}
Assumption environment: test condition setting.
\end{assumption}

% 15. Property
\begin{property}
Property environment: test object characteristics.
\end{property}

% 16. Solution
\begin{solution}
Solution to Exercise \ref{exer:test}: The sum is $55$ (test answer).
\end{solution}

% 17. Proof
\begin{proof}
Proof of Theorem \ref{thm:test}: Since $x>0$, add 1 to both sides to prove it.
\end{proof}

% ========== Chapter 4: Cross-references, Figures & Marginal Notes Testing ==========
\chapter{Cross-references and Extended Functions}
\begin{introduction}
\item All Types of Cross-references
\item Simple Figure and Table Test
\item Marginal Note Function Test
\item Chapter Jump Verification
\end{introduction}

\section{All Cross-references Summary}
Definition \ref{def:test}, Theorem \ref{thm:test}, Lemma \ref{lem:test}, Corollary \ref{cor:test}, Axiom \ref{axi:test}, Postulate \ref{pos:test}, Proposition \ref{pro:test}, Example \ref{exa:test}, Problem \ref{probl:test}, Exercise \ref{exer:test}.

\section{Figure and Table Testing}
\begin{figure}[h]
    \centering
    \includegraphics[width=0.5\textwidth]{../../image/scatter.jpg}
    \caption{Test Figure}\label{fig:test}
\end{figure}
\begin{table}[h]
    \centering
    \begin{tabular}{c*{6}{c}}
    \toprule
        $n$ & 1 & 2 & 3 & 4 & 5 & 6 \\
    \midrule
        $a_n$ & 1 & 1 & 2 & 3 & 5 & 8 \\ 
    \bottomrule
    \end{tabular}
    \caption{Test Table}\label{tab:test}
\end{table}
Figure and table reference: \figref{fig:test} is the test figure, \tabref{tab:test} is the test table.

% ========== End-of-chapter Exercises ==========
\begin{problemset}[Test Exercises of This Chapter]
\item Verify the display style of Definition \ref{def:test}
\item Derive Formula \eqref{eq:double}
\item Check if Figure \ref{fig:test} is displayed normally
\end{problemset}

% ========== References ==========
\nocite{*}
\printbibliography[heading=bibintoc]

% ========== Appendix ==========
\appendix
\chapter{Test Appendix}
The appendix content is used to verify appendix chapters, headers, page numbers, and typesetting styles, with no actual information.
\section{Appendix Test 1}
The appendix content is used to verify appendix chapters, headers, page numbers, and typesetting styles, with no actual information.
\newpage
\section{Appendix Test 2}
The appendix content is used to verify appendix chapters, headers, page numbers, and typesetting styles, with no actual information.
\end{document}